\centerline{\bf \Huge Вспоминаем радикальные оси}

\begin{flushright}
        \scriptsize{Что угодно только не ТЧ}
\end{flushright}

\problem{1}
В треугольнике $ABC$ провели высоты $AA_1$, $BB_1$ и $CC_1$, точка $H$ --- ортоцентр.
\begin{itemize}
    \item[(a)] Пусть $A'$ --- точка пересечения $B_1C_1$ и $BC$. Точки $B'$ и $C'$ определяются аналогично. Докажите, что точки $A'$, $B'$ и $C'$ лежат на одной прямой.
    \item[(b)] На сторонах $BC$, $CA$ и $AB$ выбраны произвольные точки $K$, $L$, $M$ соответственно. Докажите, что общие хорды окружностей, построенных на отрезках $AK$, $BL$ и $CM$ как на диаметрах, пересекаются в точке $H$.
\end{itemize}

\problem{2}
Для треугольника $ABC$ обозначим через $A_1$, $B_1$, $C_1$ основания высот. Пусть $A'$ --- точка пересечения касательных к окружности $(ABC)$ в точке $A$ и к окружности $(A_1B_1C_1)$ в точке $A_1$. Точки $B'$ и $C'$ определяются аналогично. Докажите, что точки $A'$, $B'$ и $C'$ лежат на одной прямой.

\problem{3}
К двум непересекающимся окружностям $\omega_1$ и $\omega_2$ проведены внешняя и внутренняя касательные $a$ и $b$, пересекающиеся в точке $C$. Обозначим точки касания окружностей $\omega_1$ и $\omega_2$ с прямой $a$ через $A_1$ и $A_2$, а с прямой $b$ --- через $B_1$ и $B_2$. Докажите, что прямые $A_1B_1$ и $A_2B_2$ пересекаются на линии центров окружностей $\omega_1$ и $\omega_2$.

\problem{4}
Вписанная окружность $\gamma$ треугольника $ABC$ касается его стороны $BC$ в точке $D$. Обозначим центр вневписанной в угол $A$ окружности через $I_a$, а середину отрезка $DI_a$ через $M$. Докажите, что отрезок касательной, проведённой из точки $M$ к $\gamma$, равен $MB$.

\problem{5}
К двум непересекающимся окружностям $\omega_1$ и $\omega_2$ проведены три общие касательные --- две внешние, $a$ и $b$, и одна внутренняя, $c$. Прямые $a$, $b$ и $c$ касаются окружности $\omega_1$ в точках $A_1$, $B_1$ и $C_1$ соответственно, а окружности $\omega_2$ --- в точках $A_2$, $B_2$ и $C_2$ соответственно. Докажите, что высоты треугольников $A_1B_1C_1$ и $A_2B_2C_2$, проведенные из вершин $C_1$ и $C_2$, равны.

\problem{6}
Дан параллелограмм $ABCD$ с тупым углом $A$. Из точки $A$ на сторону $BC$ опущен перпендикуляр $AH$. Далее, в треугольнике $ABC$ из точки $C$ проведена медиана, продолжение которой пересекает окружность $(ABC)$ в точке $K$. Докажите, что точки $K$, $H$, $C$ и $D$ лежат на одной окружности.


\newpage

\hspace{3mm}

\problem{7}
Пусть $ABC$ --- неравнобедренный треугольник с ортоцентром $H$ и описанной окружностью $\Omega$. Прямая, проходящая через $H$, пересекает отрезки $AB$ и $AC$ в точках $E$ и $F$ соответственно. Пусть $K$ --- центр окружности $(AEF)$, и пусть прямая $AK$ вторично пересекает $\Omega$ в точке $D$. Докажите, что прямая $HK$ и прямая, проходящая через $D$ перпендикулярно $BC$, пересекаются на $\Omega.$